% !TeX root = ../main.tex

\chapter{Introducción} 
\label{cap:introduccion}

    Los cúmulos estelares son conjuntos de estrellas coetáneas, unidas por su 
    propia gravedad tras formarse a partir del colapso de una misma nube de gas 
    y polvo. Al compartir edad, composición química y distancia, constituyen 
    poblaciones prácticamente libres de la degeneración que afecta a las 
    estrellas de campo, lo que los convierte en laboratorios naturales para 
    poner a prueba modelos de estructura y evolución dinámica. Su distribución 
    espacial, cinemática y edades codifican además información sobre el 
    ensamblaje y la historia química de la Vía Láctea, por lo que su 
    caracterización sistemática constituye una herramienta central de la 
    arqueología galáctica. El Capítulo \ref{cap:marco_teorico} desarrolla en 
    detalle los fundamentos físicos de estos sistemas y de los modelos 
    empleados para describirlos; basta aquí con señalar que su estudio depende 
    críticamente de la calidad de la astrometría disponible.

    En este sentido, la misión Gaia de la ESA transformó el panorama 
    observacional al proveer posiciones, paralajes y movimientos propios de una 
    precisión sin precedentes para cientos de millones de estrellas 
    \cite{Brown2018}. Combinada con algoritmos de agrupamiento no supervisado, 
    esta astrometría permitió construir catálogos homogéneos de cúmulos a una 
    escala impensable con misiones anteriores. El más reciente de estos 
    esfuerzos, el catálogo de Hunt y Reffert (2023) \cite{hunt2023improving} a 
    partir de Gaia DR3, reconoce 7167 cúmulos y reúne 1\,291\,929 estrellas 
    miembro, consolidándose como una de las bases de datos más completas 
    disponibles para el estudio estructural y dinámico de cúmulos estelares en 
    la Galaxia.

    No obstante, catálogos de esta escala privilegian necesariamente la 
    homogeneidad y la automatización sobre la caracterización dinámica 
    individual de cada sistema. Los parámetros estructurales reportados por 
    estos catálogos, radios de núcleo, radios de marea, parámetros de 
    concentración, se derivan típicamente del ajuste de perfiles de densidad 
    proyectados sobre el plano del cielo, sin una validación sistemática de si 
    dicha estructura resulta también consistente con la cinemática interna 
    observada. Asimismo, es poco frecuente que un mismo procedimiento se aplique 
    de forma homogénea tanto a cúmulos abiertos como a cúmulos globulares, pese 
    a que ambas poblaciones comparten el mismo formalismo físico de King y de 
    Jeans; esta separación metodológica dificulta comparar directamente su 
    estructura y su estado dinámico bajo un marco común. Finalmente, cuando se 
    buscan relaciones entre la estructura de los cúmulos y su entorno galáctico, 
    el tipo morfológico, estrechamente asociado a la masa, la edad y la posición 
    en la Galaxia, puede actuar como variable de confusión si no se controla 
    explícitamente por él, generando correlaciones aparentes que no reflejan 
    necesariamente un efecto ambiental real.

    Este trabajo aborda estos tres vacíos mediante un método que ajusta el 
    perfil de densidad de King \cite{king1962structure} tanto en su forma 
    proyectada (2D) como en una reconstrucción tridimensional obtenida por 
    proyección gnomónica y muestreo de rechazo de Monte Carlo (hasta 10\,000 
    realizaciones por cúmulo), aceptando como válidas únicamente las soluciones 
    que satisfacen simultáneamente criterios de bondad de ajuste ($\chi^2_\nu 
    \in [0{,}3,\,3{,}0]$, $p\geq0{,}05$) y de consistencia física entre los 
    radios de marea 2D y 3D ($r_{t,\, 3{\rm D}}\leq 5\,r_{t,\, 2{\rm D}}$). La 
    estructura así obtenida se contrasta de forma independiente con un 
    modelamiento cinemático basado en la ecuación de Jeans esférica isotrópica 
    \cite{Jeans1919}, Ecuacion~\eqref{eq:jeans_integral}, y se complementa con 
    el cálculo del tiempo de relajación de dos cuerpos \cite{Spitzer1987} para 
    evaluar el estado dinámico de cada sistema. El procedimiento se aplica bajo 
    un mismo marco metodológico a una muestra combinada de cúmulos abiertos y 
    globulares del catálogo de Hunt y Reffert, controlando explícitamente por 
    tipo morfológico en el análisis de correlaciones estructurales, evolutivas y 
    ambientales.

    \section{Objetivo General}

        Caracterizar estructural y dinámicamente una muestra de cúmulos 
        estelares abiertos y globulares de la Vía Láctea a partir de datos 
        astrométricos de Gaia DR3 (catálogo de Hunt y Reffert, 2023), mediante 
        el ajuste de perfiles de densidad de King en dos y tres dimensiones y su 
        validación cinemática con la ecuación de Jeans isotrópica, evaluando su 
        relación con la edad, el tiempo de relajación de dos cuerpos y el 
        entorno galáctico.

    \section{Objetivos Específicos}

        \begin{enumerate}
            \item Construir una muestra depurada de cúmulos estelares a partir 
            del catálogo de Hunt y Reffert (2023), aplicando criterios de 
            riqueza estelar y simetría morfológica.
            \item Implementar un método de ajuste del perfil de densidad de 
            King, en su forma proyectada (2D) y en su reconstrucción 
            tridimensional mediante proyección gnomónica y muestreo de rechazo 
            de Monte Carlo.
            \item Validar los radios estructurales, el parámetro de 
            concentración y, para los cúmulos abiertos, las masas totales 
            obtenidas, frente a los valores reportados en el catálogo de origen.
            \item Modelar el perfil de dispersión de velocidades mediante la 
            ecuación de Jeans isotrópica y contrastar los radios cinemáticos 
            ($r_{c,\, \sigma^2}$, $r_{t,\, \sigma^2}$) con los radios 
            estructurales del ajuste 3D.
            \item Determinar el tiempo de relajación de dos cuerpos 
            ($t_{\rm relax}$) y el radio de relajación ($r_{\rm relax}$) de cada 
            sistema para caracterizar su estado dinámico.
            \item Analizar las relaciones entre los parámetros estructurales y 
            dinámicos y las propiedades físicas y ambientales de los cúmulos 
            (edad, masa, radio galactocéntrico, altura y latitud galáctica), 
            controlando por tipo morfológico.
        \end{enumerate}

    \section{Alcance y Limitaciones}

        El análisis se restringe a cúmulos con al menos 1000 miembros 
        confirmados y con simetría espacial suficiente para justificar un ajuste 
        esférico, criterios que redujeron el catálogo original a una muestra 
        final de 100 cúmulos (54 abiertos y 46 globulares). El ajuste 
        tridimensional y su validación cinemática, al depender de la 
        convergencia de un procedimiento estocástico bajo criterios de calidad, 
        solo resultaron estadísticamente aceptados para 59 de estos 100 
        sistemas, por lo que las conclusiones de naturaleza dinámica se apoyan 
        principalmente en esta submuestra. El modelamiento cinemático asume, 
        además, simetría esférica e isotropía en la distribución de velocidades, 
        y la reconstrucción de la profundidad tridimensional depende de las 
        distancias fotométricas bayesianas de Bailer-Jones et al. (2018) 
        \cite{bailer2018estimating}, cuyas incertidumbres se propagan a los 
        parámetros estructurales derivados. Finalmente, la ausencia de edades 
        fotométricas para los cúmulos globulares en el catálogo de origen 
        restringe el análisis evolutivo, la relación entre estructura, 
        relajación dinámica y edad, exclusivamente a la submuestra de cúmulos 
        abiertos.

    \section{Estructura del Documento}

        El resto del documento se organiza de la siguiente manera. El 
        Capítulo~\ref{cap:marco_teorico} presenta el marco teórico: la misión 
        Gaia, las propiedades físicas de los cúmulos estelares, el modelo de 
        densidad de King y la formulación de las ecuaciones de Jeans para 
        sistemas esféricos. El Capítulo~\ref{cap:metodologia} describe el 
        catálogo de Hunt y Reffert, los criterios de selección y limpieza de la 
        muestra, y el método de procesamiento: la proyección gnomónica, el 
        ajuste del perfil de King en dos y tres dimensiones mediante muestreo de 
        rechazo de Monte Carlo, y el modelamiento cinemático a través de la 
        ecuación de Jeans. El Capítulo~\ref{cap:resultados} presenta los 
        resultados: la caracterización general de la muestra, la validación del 
        método frente al catálogo, las relaciones de estructura interna, las 
        correlaciones evolutivas y ambientales, y los tiempos de relajación de 
        dos cuerpos, cerrando con una síntesis integradora de los hallazgos. 
        Finalmente, el Capítulo~\ref{cap:conclusiones} sintetiza las 
        conclusiones del trabajo, discute sus limitaciones y propone líneas de 
        trabajo futuro.